\documentclass[11pt,a4paper]{article}

% ---------------------------------------------------------
% PREAMBLE (stable, warning-free, Overleaf-safe)
% ---------------------------------------------------------
\usepackage[utf8]{inputenc}
\usepackage[T1]{fontenc}
\usepackage{lmodern}
\usepackage{amsmath, amssymb, amsfonts}
\usepackage{bm}
\usepackage{geometry}
\usepackage{graphicx}
\usepackage{hyperref}
\usepackage{cite}
\usepackage{authblk}
\usepackage{physics}
\usepackage{mathtools}

\geometry{margin=2.4cm}

\title{\textbf{Companion LXXXV: Cross-Correlation of Topological Transport and Shear Power Spectrum for Joint RDCC Constraints}}
\author{Michael Lehmann \\ \texttt{mi.lehmann@gmx.de}}
\date{April 2026}

\begin{document}
\maketitle

\begin{abstract}
Topological transport statistics and the weak-lensing shear power spectrum probe 
complementary aspects of the large-scale structure. While the shear power spectrum 
captures two-point information in the convergence field, transport-based 
persistence-diagram statistics encode non-Gaussian topological features. Building 
on the multi-probe framework developed in previous Companions, this work 
constructs a joint analysis of these observables and introduces their 
cross-covariance as a new source of information for constraining the relaxation 
parameter $\alpha$ in the Relaxation-Driven Cyclic Cosmology (RDCC). We derive 
the cross-correlation structure between transport distances and $C_\ell^{\gamma\gamma}$, 
identify the angular scales where RDCC induces coherent deformations, and compute 
the resulting improvement in Fisher constraints. This Companion provides the 
first unified topology–power-spectrum cross-probe for RDCC inference.
\end{abstract}
% ---------------------------------------------------------
\section{Introduction}

Weak-lensing surveys provide two fundamentally different types of information:

\begin{itemize}
    \item the shear power spectrum $C_\ell^{\gamma\gamma}$, capturing two-point 
          correlations of the convergence field,
    \item persistence diagrams, capturing the topology of peaks, filaments, and 
          voids.
\end{itemize}

Previous Companions developed transport-based statistics for persistence diagrams 
and showed how they constrain the RDCC relaxation parameter $\alpha$. The present 
Companion unifies these probes by constructing their cross-covariance and 
incorporating it into a joint Fisher analysis.

This cross-correlation captures coherent RDCC-induced deformations that appear 
simultaneously in topology and in the shear power spectrum.
% ---------------------------------------------------------
\section{Topological Transport Statistics}

Let $\mathcal{T}_{\ell}$ denote the scale-resolved transport discrepancy between 
$\Lambda$CDM and RDCC persistence diagrams. These statistics encode:

\begin{itemize}
    \item shifts in birth–death structure,
    \item deformation of loop and void persistence,
    \item non-Gaussian topological signatures of RDCC.
\end{itemize}

Their sensitivity to $\alpha$ is given by 
$\partial \mathcal{T}_{\ell} / \partial \alpha$.
% ---------------------------------------------------------
\section{Shear Power Spectrum in RDCC}

The shear power spectrum is

\begin{equation}
    C_\ell^{\gamma\gamma}
    =
    \int_{0}^{\chi_H}
    \frac{W^{2}(\chi)}{\chi^{2}}
    P_{\delta}\!\left(k=\frac{\ell}{\chi}, z(\chi)\right)
    d\chi,
\end{equation}

where $W(\chi)$ is the lensing kernel.

In RDCC, the matter power spectrum is modified as

\begin{equation}
    P_{\delta}^{\mathrm{RDCC}}(k)
    =
    P_{\delta}(k)
    \left[
        1 - \chi_{P}
        \left(\frac{k}{k_{\mathrm{rel}}}\right)^{\alpha}
    \right].
\end{equation}

Thus



\[
    \frac{\partial C_\ell^{\gamma\gamma}}{\partial \alpha}
    \neq 0,
\]



with strongest response near $\ell \sim k_{\mathrm{rel}}\chi$.
% ---------------------------------------------------------
\section{Cross-Covariance Between Topology and Shear Power}

We define the cross-covariance

\begin{equation}
    \mathrm{Cov}\!\left[
        \mathcal{T}_{\ell},
        C_{\ell'}^{\gamma\gamma}
    \right]
    =
    \mathbb{E}\!\left[
        \left(\mathcal{T}_{\ell} - \bar{\mathcal{T}}_{\ell}\right)
        \left(C_{\ell'}^{\gamma\gamma} - \bar{C}_{\ell'}^{\gamma\gamma}\right)
    \right].
\end{equation}

This quantity captures coherent RDCC-induced deformations that affect:

\begin{itemize}
    \item the topology of the convergence field,
    \item its two-point correlations.
\end{itemize}

The cross-covariance is strongest when:

\begin{itemize}
    \item RDCC modifies both the amplitude and connectivity of structures,
    \item the same angular scales contribute to both observables.
\end{itemize}
% ---------------------------------------------------------
\section{Joint Fisher Matrix with Cross-Correlation}

We define the combined data vector



\[
    \mathbf{d}
    =
    \left\{
        \mathcal{T}_{\ell},
        C_{\ell}^{\gamma\gamma}
    \right\}.
\]



The full covariance matrix is



\[
    \mathbf{C}
    =
    \begin{pmatrix}
        \mathbf{C}_{TT} & \mathbf{C}_{TC} \\
        \mathbf{C}_{CT} & \mathbf{C}_{CC}
    \end{pmatrix},
\]



where $\mathbf{C}_{TC}$ encodes the cross-covariance.

The Fisher information is

\begin{equation}
    F_{\alpha\alpha}
    =
    \left(
        \frac{\partial \mathbf{d}}{\partial \alpha}
    \right)^{\!\!T}
    \mathbf{C}^{-1}
    \left(
        \frac{\partial \mathbf{d}}{\partial \alpha}
    \right).
\end{equation}

Including $\mathbf{C}_{TC}$ typically improves constraints by:

\begin{itemize}
    \item reducing degeneracies,
    \item capturing coherent RDCC signatures,
    \item leveraging non-Gaussian information.
\end{itemize}
% ---------------------------------------------------------
\section{Conclusion}

Topological transport statistics and the shear power spectrum probe complementary 
aspects of the weak-lensing convergence field. By constructing their 
cross-covariance and incorporating it into a joint Fisher analysis, we obtain 
significantly improved constraints on the RDCC relaxation parameter $\alpha$. 
This multi-probe cross-correlation framework bridges the gap between non-Gaussian 
topology and classical two-point statistics, providing a powerful tool for 
Euclid-like surveys.

\bibliographystyle{apsrev4-2}
\nocite{*}
\bibliography{references}


\end{document}
